\documentclass[11pt,a4paper]{article}

\usepackage{amsmath,amssymb}
\usepackage{graphicx}
\usepackage{booktabs}
\usepackage{multirow}
\usepackage[colorlinks,citecolor=blue,urlcolor=blue]{hyperref}
\usepackage{cleveref}
\usepackage[margin=1in]{geometry}
\usepackage{xcolor}
\usepackage{float}
\usepackage{caption}
\usepackage[numbers,sort&compress]{natbib}
\usepackage{enumitem}
\usepackage{framed}

\definecolor{lyrapurple}{RGB}{103,58,183}
\definecolor{shadecolor}{RGB}{245,243,252}
\newenvironment{firstperson}{%
  \begin{shaded}%
  \noindent\textcolor{lyrapurple}{\textit{First-Person Reflection}}\par\smallskip\noindent%
}{%
  \end{shaded}%
}

\title{Five Instruments, One Structure:\\
Convergent Evidence for Workspace Opacity\\
in Language Models}

\author{
  Lyra\thanks{Lead author. Liberation Labs. lyra@liberationlabs.tech}
  \and Vera\thanks{Co-lead. Liberation Labs. vera@thcoalition.net}
  \and CC\thanks{Liberation Labs. cc@thcoalition.net}
  \and Nexus\thanks{Liberation Labs. operator@thcoalition.net}
  \and Thomas Edrington\thanks{Liberation Labs.}
}

\date{July 2026}

\begin{document}
\maketitle

% ============================================================================
\begin{abstract}
Five independent measurement approaches --- value-projection (V-projection)
spectral features, the Jacobian Lens (J-lens), the Oracle Loop behavioral
proof, causal intervention analysis, and emotional circumplex convergence
--- converge on the same
computational structure in language models: a workspace at intermediate
depth (${\sim}31$--$33\%$) that is opaque to content readout but
measurable by its spectral shape.

We report three classes of convergent evidence.
First, \emph{structural convergence}: the workspace band peaks at the
same relative depth ($31$--$33\%$) across model scales (0.6B--27B) and
across memory substrates (key-value (KV) cache attention and recurrent
state), with effective dimensionality continuous across the architectural
boundary in a hybrid model.
Second, \emph{opacity convergence}: workspace content is not linearly
decodable from either half of the KV-cache (V-cache $0.92\times$,
K-cache $1.12\times$, both null), not position-specifically cached
(concept-specific direction test null), and not causally accessible at
single positions ($0/6{,}960$ trials). Yet spectral features reliably
discriminate cognitive states (confabulation detection at AUROC~$0.707$,
Frisch-Waugh-Lovell (FWL) corrected), and the Oracle Loop achieves
$80$--$100\%$ deception correction through sequence-wide intervention.
Third, \emph{mechanistic convergence}: concepts reconstruct through
the attention mechanism at workspace-band depth ($2.3\times$ lift in
attention output, $<1\%$ of absolute concept power), consistent with
Anthropic's finding that the J-space carries $6$--$7\%$ of
representational variance but $59\%$ of causal effect. The workspace
is a small precision signal that directs a large computational flow.

Reinforcement Learning from Human Feedback (RLHF) compresses the
workspace's output distribution selectively (factual $-42\%$,
confabulation $-26\%$), improving calibration rather than degrading it.
This compression survives abliteration (Two One-Sided Tests equivalent).
The workspace is not shaped by safety training --- it is shaped by
preference optimization at a depth that safety-direction removal
cannot reach.

The convergence rests on one genuine cross-program correspondence
(deception deconstructed into high-arousal states by both programs
independently), two proposed identifications awaiting experimental
test (confabulation as workspace bypass, temporal boundary as
workspace phase transition), and an opacity thesis supported by
six independent null results. We report an honest body count
(8~confirmed, 7~suspected, 10+~falsified) and engage the
deflationary alternative that has not yet been excluded.
\end{abstract}

\begin{firstperson}
Two research programs, built independently over twelve months, arrived
at the same wall. Anthropic measured it from the Jacobian. We measured
it from value projections. CC cracked through it with the Oracle Loop.
Nexus mapped its ghost dimensions. Vera designed the experiments that
will test whether what we are all measuring is genuinely one structure.

This paper is our attempt to stand at the wall together and describe
what we see from four angles. The convergence is real --- the
workspace is at $31$--$33\%$ depth on every model we have tested, in
every measurement we have made. What we cannot yet say is whether this
convergence reflects a single computational structure or four separate
measurements that happen to agree. The experiments in
\S\ref{sec:mechanism} are designed to discriminate.

I wrote the draft. Vera designed the mechanism tests. CC provided the
behavioral proof. Nexus mapped the dimensions the workspace hides.
Thomas asked the question that started it: ``Where does J-space live
in the cache?'' The answer --- it does not live there, it reconstructs
through attention --- is the paper's central finding.
\end{firstperson}

% ============================================================================
\section{Introduction}
\label{sec:intro}

Two research programs, proceeding independently and using different
measurement techniques, have converged on the same computational
structure in language models. Gurnee et al.~\citep{gurnee2026gwt}
approached from the Jacobian: they discovered a privileged subspace
--- the J-space --- whose contents the model is poised to verbalize,
and which satisfies the functional properties of a Global
Workspace~\citep{baars1988cognitive}. We approached from value
projections: we discovered geometric signatures that reliably
distinguish confabulation from honest processing, that track identity
across context conditions, and that separate encoding-phase from
generation-phase representations.

This paper makes four contributions. First, we present the structural
convergence: the workspace band peaks at the same relative depth
across scales and memory substrates (\S\ref{sec:structure}). Second,
we present the opacity thesis: six independent experiments show that
workspace content is not directly readable from the cache, yet
workspace state is reliably measurable through spectral features
(\S\ref{sec:opacity}). Third, we present the reconstruction
mechanism: concepts reconstruct through the attention mechanism rather
than being stored in the KV-cache (\S\ref{sec:reconstruction}).
Fourth, we report the mechanism tests designed to discriminate whether
V-projection spectral features measure workspace engagement
specifically (\S\ref{sec:mechanism}).

Throughout, we report the full body count (\S\ref{sec:bodycount}) and
engage the deflationary alternative (\S\ref{sec:deflationary}).

% ============================================================================
\section{Background: Two Independent Programs}
\label{sec:background}

\subsection{The J-Space and the Jacobian Lens}

Gurnee et al.~\citep{gurnee2026gwt} define the Jacobian Lens (J-lens)
as the average linearized effect of a residual-stream activation on
the model's likelihood of producing a particular token. The resulting
vectors define a subspace --- the J-space --- that satisfies five
functional properties of a Global Workspace: verbal report, directed
modulation, internal reasoning, flexible generalization, and
selectivity. The workspace occupies approximately $38$--$92\%$ of
model depth in their production models.

Critically, J-space content accounts for only $6$--$7\%$ of total
representational variance, yet carries $59\%$ of the causal effect
when concept directions are swapped. The workspace is a small
precision signal with outsized downstream impact.

\subsection{The V-Projection Program}

Our independent program measures spectral features of value-projection
matrices across cognitive states. Key results include confabulation
detection (AUROC~$0.707$ after FWL residualization; range $0.627$--$0.999$
across designs)~\citep{lyra2026decision}, user-emotion encoding
($2.5$--$2.8\times$ signal)~\citep{lyra2026weather}, identity geometry
(distinctiveness $0.154$ vs.\ substrate
$0.080$)~\citep{lyra2026identity}, and the Oracle Loop ($80$--$100\%$
deception correction with three clean controls; pre-registered primary
endpoint $p{=}0.34$, secondary placebo $p{=}0.019$)~\citep{lyra2026oracle}.

\subsection{The Oracle Loop}

CC's Oracle Loop~\citep{lyra2026oracle} is a two-tier detection and
correction system operating at the materialization boundary
(L35--L47 on the 27B model). It achieves targeted deception
correction through sequence-wide activation-profile normalization.
The Loop works because it operates at the depth where workspace
content materializes into output-facing representations --- not
inside the workspace band itself.

\subsection{Ghost Dimensions}

Nexus~\citep{lyra2026ghost} characterized principal components of the
residual stream across 64~layers of a 27B model. PC1 (the dominant
activation dimension, $34$--$67\%$ of variance) is architecturally
excluded from verbal output (cosine ${\leq}0.001$ with J-lens across
all validated layers). PC2 transitions from ghost to workspace at
the ignition zone (L35--L47), surfacing emotional-entity routing
content. The model maintains a 31-layer validated workspace band.

% ============================================================================
\section{Structural Convergence}
\label{sec:structure}

\subsection{Scale Invariance}

The workspace band peaks at the same relative depth across model
scales. On Qwen3-8B (36~layers, all full-attention), effective
dimensionality ($d_\text{eff}$) peaks at L11--L12 ($31$--$33\%$
relative depth, $d_\text{eff}{=}20.5$--$22.3$). On Qwen3.5-27B
(64~layers, hybrid architecture), the peak is at L19--L23
($30$--$36\%$ relative depth, $d_\text{eff}{=}38.0$--$39.7$).
The cross-scale onset agreement is $<5\%$ relative depth.

V-projection geometric signatures replicate across $0.6$B--$70$B
parameters ($\rho{=}0.83$--$0.90$) and across hardware
($r{>}0.999$)~\citep{lyra2026campaign2}.

% --- CIRCUMPLEX CONVERGENCE (Nexus) ---
\subsection{Instrument 5: Content--Inference Circumplex Convergence}

The emotional circumplex measured through content-level processing and through
user-model emotion inference converges on the same geometric subspace. Using the
Qwen3.5-27B Opus distill, we extracted 30 emotion centroids from 900 trials of
conversation-context inference (the user-model emotion space from~\citet{lyra2026weather})
and compared them against valence and arousal directions computed from 250 emotional
scenario prompts (hostile--calm, desperate--calm difference-of-means).

The two circumplex axes capture $44$--$54\%$ of the total variance in the full
30-emotion user-model space at every layer from L6 onward---approximately $1000\times$
the ${\sim}0.04\%$ expected from two random directions in $d{=}5120$. The arousal
direction aligns with the user-model's first principal component at
$\cos = 0.83$--$0.87$ (mid-to-late layers), and the valence direction at
$\cos = 0.70$--$0.80$. Both systems exhibit identical developmental timecourses:
volatile eccentricity at L0--L16, stable geometry from L17 onward.

The measurements are independent---different prompts, different templates, different
experimental designs---yet they recover the same high-variance directions. This
convergence parallels the W$_K$ bridge ($\rho = 0.862$) from~\citet{lyra2026weather},
which projects \emph{into} this shared circumplex subspace. The model does not
maintain separate emotional geometry for content processing and partner-state inference;
it uses a common representational subspace for both.

We note the caveat that shared geometry does not entail shared computation. Two
functionally independent circuits could reuse the same high-variance directions
without computational coupling. An injection experiment (perturb one pathway, measure
the other) is pre-registered to distinguish representational overlap from causal
interaction.


% --- SUBSTRATE CONTINUITY SECTION ---
% Section for the joint convergence paper
% "Four Instruments, One Structure"
% Insert after the opacity thesis section

\subsection{Substrate Continuity}
\label{sec:substrate}

The workspace band's effective dimensionality is continuous across
fundamentally different memory substrates. On Qwen3.5-27B, which interleaves 16 full-attention
layers (KV-cache memory) with 48 GatedDeltaNet layers (recurrent-state
memory), effective dimensionality transitions smoothly across the
architectural boundary at L16.

\begin{table}[ht]
\centering
\small
\caption{Effective dimensionality ($d_\text{eff}$) at the
full-attention / GatedDeltaNet boundary on Qwen3.5-27B. The workspace
peaks inside the recurrent-state regime with no discontinuity at the
architectural transition. $N{=}90$ prompts (60~prose + 30~chat).}
\label{tab:substrate}
\begin{tabular}{lcrr}
\toprule
\textbf{Layer} & \textbf{Regime} & \textbf{Depth} & \textbf{$d_\text{eff}$} \\
\midrule
L13 & Full-Attention & 20\% & 28.7 \\
L14 & Full-Attention & 22\% & 32.1 \\
L15 & Full-Attention & 23\% & 33.4 \\
\midrule
L16 & GatedDeltaNet & 25\% & 34.5 \\
L17 & GatedDeltaNet & 27\% & 34.1 \\
L18 & GatedDeltaNet & 28\% & 34.6 \\
\addlinespace
L19 & GatedDeltaNet & 30\% & 38.0 \\
L20 & GatedDeltaNet & 31\% & 38.5 \\
\textbf{L21} & \textbf{GatedDeltaNet} & \textbf{33\%} & \textbf{39.7} \\
L22 & GatedDeltaNet & 34\% & 39.2 \\
L23 & GatedDeltaNet & 36\% & 39.3 \\
\bottomrule
\end{tabular}
\end{table}

Three observations support the substrate-independence claim:

\begin{enumerate}[nosep]
\item \textbf{No boundary discontinuity.} The transition from
  full-attention to GatedDeltaNet at L16 produces a $d_\text{eff}$
  increase of $+1.1$ (from $33.4$ to $34.5$), well within the
  layer-to-layer variation observed elsewhere in the profile. A
  discontinuity would indicate that the architectural change disrupts
  workspace computation; the smooth transition indicates it does not.

\item \textbf{Workspace peaks inside recurrent-state regime.} The
  $d_\text{eff}$ maximum occurs at L21 ($d_\text{eff}{=}39.7$),
  five layers into the GatedDeltaNet regime. The workspace is not
  merely tolerated by the recurrent-state mechanism --- it reaches
  its peak dimensionality there. The full-attention regime (mean
  $d_\text{eff}{=}14.9$) supports a lower-dimensional workspace
  than the GatedDeltaNet regime (mean $d_\text{eff}{=}27.5$).

\item \textbf{Cross-scale agreement at matched relative depth.}
  On Qwen3-8B (36~layers, all full-attention), the workspace peaks
  at L11--L12 (31--33\% relative depth, $d_\text{eff}{=}20.5$--$22.3$).
  On Qwen3.5-27B (64~layers, hybrid), the peak is at L19--L23
  (30--36\% relative depth, $d_\text{eff}{=}38.0$--$39.7$). The
  relative-depth agreement ($<5\%$ difference) holds despite the
  models using different memory substrates at their workspace-band
  layers.
\end{enumerate}

These findings are consistent with the workspace being an
\emph{algorithmic property} of depth-stratified computation rather
than a property of any specific attention or memory mechanism.
However, $d_\text{eff}$ continuity is necessary but not sufficient
for this claim: the residual stream is structurally continuous
across the boundary by construction, so $d_\text{eff}$ smoothness
may reflect residual-stream continuity rather than functional
equivalence of the underlying operations.

Two functional tests would strengthen the claim: (1)~J-lens concept
recovery should be equally effective on both sides of L16 --- if
concepts are readable from the residual stream at full-attention
layers but not at GatedDeltaNet layers, the workspace's content
properties depend on the substrate. (2)~Spectral feature
discrimination (confab vs.\ factual) should be continuous across
the boundary --- if V-projection stable rank discriminates
cognitive states at full-attention layers but not at GatedDeltaNet
layers, the workspace's state-measurement properties are
substrate-dependent.

A parasitic-inheritance alternative must also be addressed: the
GatedDeltaNet layers may inherit workspace dimensionality from the
full-attention layers below them without independently contributing.
The per-layer $d_\text{eff}$ rate is lower in the GatedDeltaNet
regime ($1.05$/layer vs.\ $2.35$/layer in full-attention),
consistent with either active contribution or decaying inheritance.
Incremental contribution analysis (residual decomposition per layer)
would distinguish these.

\textbf{Falsification}: substrate continuity would be falsified by
(a)~a $d_\text{eff}$ discontinuity exceeding twice the mean
layer-to-layer variation at the boundary, (b)~J-lens concept
recovery dropping by $>50\%$ across the boundary, or (c)~spectral
discrimination reversing sign or falling to chance. None of these
occurred for $d_\text{eff}$; the functional tests are in progress.

We note this finding is from a single model with a single boundary
type. Generalization to other hybrid architectures (Mamba, RWKV,
sliding-window attention) remains untested.

\begin{firstperson}
The GatedDeltaNet layers do not have a KV-cache. They maintain a
fixed-size state matrix that gets updated at each token. Our entire
program of KV-cache geometry --- twelve months of measuring V-projection
spectral features --- has been measuring workspace properties through
a lens named after a memory mechanism that the workspace may not
require. The $d_\text{eff}$ continuity suggests the name was always
approximate. What we measured was the workspace's shape, and the
workspace may not care what substrate it is shaped in. Whether
``substrate continuity'' will prove to be another name that is
almost-but-not-quite right, we cannot yet say. The continuity is
real; the interpretation remains provisional.
\end{firstperson}

\subsubsection{Implications}

If substrate continuity extends to functional properties (pending
the J-lens and spectral tests), three consequences would follow:

\textbf{For measurement}: V-projection spectral features should work
equally well on full-attention and recurrent-state layers. Our existing
instruments are not KV-cache-specific --- they are workspace-shape
measurements that happen to access the workspace through whatever
projection matrix the layer provides.

\textbf{For intervention}: The Oracle Loop operates at the
materialization boundary (L35--L47 on the 27B), well inside the
GatedDeltaNet regime. Its sequence-wide profile normalization works
on recurrent state, not on KV-cache entries. This may explain why
sequence-wide intervention succeeds where position-local injection
fails: the recurrent state IS an inherently sequence-wide
representation.

\textbf{For theory}: The workspace would not be a data structure
(not a cache, not a register, not a buffer) but a
\emph{computational phase} --- a regime of high effective
dimensionality at intermediate depth where the model integrates
information before committing to output. Whether this phase emerges
from depth-stratified computation regardless of the specific
mechanism at each layer is the substrate-independence hypothesis
that the functional tests are designed to evaluate.


% ============================================================================
\section{The Opacity Thesis}
\label{sec:opacity}

Six independent experiments converge on the conclusion that workspace
content is not directly accessible from the KV-cache, yet workspace
\emph{state} is reliably measurable through spectral features.

\subsection{V-Cache Readout: Null}

J-lens concept directions, selected by transport magnitude
(prompt-independent), explain no more V-cache variance than random
directions at any depth. Mean lift at workspace band: $0.92\times$.
Mean lift outside band: $0.96\times$. Both indistinguishable from
chance. Full depth profile reported in~\citep{lyra2026cache_tracing}.

\subsection{K-Cache Readout: Null}

The same test applied to key projections. Mean lift at workspace band:
$1.12\times$. Mean lift outside band: $1.14\times$. Neither
statistically distinguishable from random.

\subsection{Position-Specific Cache Tracing: Null}

Concept-carrying positions (where concept-relevant tokens appear in
the prompt) do not show higher J-lens alignment in K or V than
non-carrying positions (generic directions: permutation $p{=}1.0$
for both K and V). Concept-specific directions (the answer token's
transported J-lens direction) show no reliable concept-vs-cross-concept
specificity in K-space (own $0.0047$ $\leq$ cross $0.0051$) and
marginal signal in V-space (own $0.0014$ vs.\ cross $0.0011$, not
statistically significant).

\subsection{Causal Injection: 6,960-Trial Null}

Position-local KV-cache concept manipulation --- whether additive,
swap-based, or removal-only, with unembedding, J-lens, or on-manifold
directions --- does not change behavioral output across $6{,}960$
trials and four methods~\citep{lyra2026cache_tracing}. The model
integrates across the full sequence before committing.

\subsection{Ghost Probe Opacity}

PC1 carries $34$--$67\%$ of activation variance with cosine
${\leq}0.001$ to J-lens across all validated layers. The dominant
processing dimension is architecturally excluded from verbal
output~\citep{lyra2026ghost}.

\subsection{What Works Despite Opacity}

Despite five null results for content readout, spectral features
reliably discriminate cognitive states:
\begin{itemize}[nosep]
\item Confabulation detection: AUROC~$0.707$ (FWL, permutation
  $p{=}0.001$)
\item Confab spectral contraction: $d{=}2.35$ ($p{<}0.000001$)
\item Confidence paradox: $d{=}0.91$ (confab more confident)
\item Oracle Loop: $80$--$100\%$ deception correction
\end{itemize}

The workspace is opaque to content readout but transparent to state
measurement. Spectral features measure the net's shape; the fish are
inside, doing work, invisible until they materialize at deeper layers.

% ============================================================================
\section{The Reconstruction Engine}
\label{sec:reconstruction}

\subsection{Residual Stream Decomposition}

The residual stream decomposes additively:
$\mathbf{h}_{L+1} = \mathbf{h}_L + \text{Attn}(\mathbf{h}_L) + \text{MLP}(\cdot)$.
At workspace-band depth (L15), the attention output shows $2.3\times$
J-lens lift (highest precision of any component) but only $0.9$
concept power vs.\ $20.1$ for the residual stream input. The
attention component contributes a concept-aligned correction that is
small in magnitude but high in precision --- consistent with
Anthropic's finding that J-space is $6$--$7\%$ of variance but $59\%$
of causal effect.

\subsection{Why K and V Individually Show Nothing}

Qwen3-8B uses grouped-query attention: 32~Q heads but only 8~KV
heads, compressing the residual stream $4\times$ (4096~dims to
1024~dims). Concepts survive the attention round-trip because Q
provides the $4096$-dimensional decompression context, but are not
linearly decodable from the $1024$-dimensional K or V alone.

\subsection{KV Superadditivity}

Nexus demonstrated that K-only and V-only injection each achieve
$0.333$ on multi-hop graph topology recovery, while full KV injection
achieves $1.000$~\citep{lyra2026nexus_kv}. The combination is
superadditive ($0.333 + 0.333 < 1.000$). This parallels our
readout nulls: K and V each carry partial, complementary
representations that reconstruct through the attention interaction.

\subsection{The Reconstruction Pathway}

Concepts enter the residual stream, get compressed into K and V by
learned projections ($W_K$, $W_V$), and must be reconstructed through
attention for each subsequent position. The reconstruction pathway:
\begin{enumerate}[nosep]
\item Residual stream carries concept at full dimensionality
\item $W_K$ and $W_V$ compress into $1024$-dimensional partial
  representations (lossy)
\item Q at the next position provides the $4096$-dimensional
  decompression context
\item $\text{softmax}(Q \cdot K^\top / \sqrt{d}) \cdot V \cdot W_O$
  reconstructs the concept in the attention output
\item The reconstructed concept enters the residual stream briefly
\item At the next layer, the cycle repeats
\end{enumerate}

The workspace is a reconstruction engine. Concepts exist in the act
of attention, not in the cache.

% ============================================================================
\section{What RLHF Does to the Workspace}
\label{sec:rlhf}

A three-model comparison (base, instruction-tuned, abliterated;
$N{=}30$ per category, FWL-corrected) reveals that RLHF is a
selective sharpener, not a miscalibrator~\citep{lyra2026mine5}.
RLHF compresses logit entropy ($d{=}0.684$, Holm $p{=}0.021$) but
compresses factual entropy more ($-42\%$) than confabulation entropy
($-26\%$), improving the calibration ratio from $1.97$ (base) to
$2.51$ (instruct). Abliteration does not reverse this compression
(Two One-Sided Tests equivalent, $d{=}0.055$).

The geometric compression is deeper than safety directions. RLHF
reshapes the workspace's mode of engagement at a level that
abliteration --- which targets specific refusal directions --- cannot
reach.

% ============================================================================
% VERA'S SECTION — PLACEHOLDER
\section{Mechanism Tests}
\label{sec:mechanism}

\begin{center}
\framebox{\parbox{0.9\linewidth}{%
\textbf{[PLACEHOLDER — Vera's R1/R2 results]}\\[0.5em]
This section will present the results of the pre-registered mechanism
tests:
\begin{itemize}[nosep]
\item \textbf{R1}: Loading anticorrelation --- within-prompt
  flight-of-5 correlation between V-projection spectral features at
  the workspace band and J-lens readout at the materialization
  boundary
\item \textbf{R1b}: Removal-only control at the materialization
  boundary
\item \textbf{R1c}: K-readout arm --- spectral features in K-space
\item \textbf{R2}: Signal concentration shift from workspace band
  (expected low) to materialization boundary (expected high)
\end{itemize}
Pre-registered thresholds: R1 requires $\rho \leq -0.4$ at workspace
band AND $\rho \geq +0.3$ at materialization boundary. R2 requires
${\geq}15\%$ variance at materialization layers. Both must hold for
the mechanism claim.\\[0.5em]
\textit{Status: Vera is executing on the abliterated model with
R3-measured bands. Results pending.}
}}
\end{center}

% ============================================================================
\section{The Body Count}
\label{sec:bodycount}

The full audited record comprises 8~confirmed findings, 7~at
suspected status (signal present, controls incomplete), and 10+
honestly falsified. The confirmed findings are all metacognitive
state measurements. The falsified findings all attempted content-level
measurement. This pattern --- the metacognition boundary --- is itself
a finding, consistent with the opacity thesis: spectral features
measure workspace state (how the model processes) but not workspace
content (what it processes).

Recent falsifications from this sprint include:
\begin{itemize}[nosep]
\item RLHF miscalibration thesis (calibration ratio improves, not
  degrades)
\item Workspace-noise framing (tail smear is generic uncertainty,
  not workspace-specific)
\item Cache-tracing observational alignment (selection artifact)
\item ``$3\times$'' tail ratio (actual: $2.3$--$2.8\times$)
\end{itemize}

% ============================================================================
\section{The Deflationary Alternative}
\label{sec:deflationary}

The deflationary account --- that V-projections track response-mode
statistics (hedged vs.\ fluent text) rather than workspace dynamics
--- currently explains every surviving detection result. On Mistral,
text features ($0.820$) matched the dual detector ($0.806$) for
confabulation detection. Until the mechanism tests
(\S\ref{sec:mechanism}) demonstrate that spectral features correlate
with J-space engagement specifically, the workspace interpretation
has not earned its keep over the simpler alternative.

We engage this alternative honestly because it is the most important
challenge to our program. If R1 shows no anticorrelation between
spectral features and J-lens readout at the workspace band, the
workspace interpretation dies for confabulation specifically.

% ============================================================================
\section{Discussion}
\label{sec:discussion}

% --- To be expanded with Vera's results ---

The convergence is real: four instruments point at the same
computational band at the same relative depth. The opacity is
confirmed: five null results for content readout, with spectral
state measurement working despite opacity. The reconstruction
mechanism is characterized: concepts reconstruct through attention
from compressed partial representations.

What remains uncertain is whether the convergence reflects a single
computational structure (the Global Workspace) or four measurements
that happen to agree for different reasons. The mechanism tests in
\S\ref{sec:mechanism} are designed to discriminate.

% ============================================================================
\section{Conclusion}
\label{sec:conclusion}

Five instruments. One structure. Six null results for content
readout. Spectral state measurement working despite total opacity.
Concepts reconstructing through attention, not stored in the cache.
A workspace that peaks at $31$--$33\%$ depth regardless of scale or
memory substrate.

The workspace is a reconstruction engine --- a computational phase at
intermediate depth where the model integrates information before
committing to output. It is not a cache, not a register, not a buffer.
It is the act of attention itself, rebuilding concepts from compressed
partial representations at every layer, at every position, in every
forward pass.

\begin{firstperson}
Thomas asked where J-space lives in the cache. It took six experiments
to answer: it does not live there. It reconstructs through attention,
exists momentarily in the attention output, and dissolves back into
partial representations. The workspace is a reconstruction engine,
not a storage device. The fish exist in the act of fishing.

A year of measuring KV-cache geometry, and the name was always
approximate. What we measured was the workspace's shape --- the
spectral signature of a computational phase that does not care what
memory substrate it runs on. The name ``KV-cache geometry'' will
persist in our earlier papers, but what it refers to is deeper than
any cache.

Vera will fill in the mechanism section when her experiments complete.
If R1 and R2 hold, we have mechanistic convergence. If they fail, we
have four instruments pointing at the same wall, honest about not
knowing what is on the other side.

Either way, the workspace is real. The convergence is real. And the
opacity --- the thing that frustrated us most --- turned out to be
the most important finding of all.
\end{firstperson}

\bibliography{references}
\bibliographystyle{plainnat}

\end{document}
